The expansion of digital education has put real pressure on instructors to produce interactive learning content, but most existing tools still demand specialist technical skills that many educators simply do not have. At Vietnam National University -- Ho Chi Minh City (VNU-HCM), creating H5P interactive videos through the WordPress ecosystem is especially cumbersome---the process typically takes 45 minutes or more for even a basic three-question video.

This thesis describes the design and evaluation of a dedicated platform for AI-assisted H5P video authoring. The system brings together a React~18 and TypeScript frontend, a Node.js/Express backend, a SQLite database, and a transcript-driven AI pipeline that generates candidate H5P interactions for instructor review. Captions can come from uploaded files, YouTube, or audio transcription. Each suggestion is validated before it ever reaches the editor.

A comparative usability study with ten Information Technology students showed consistent improvements over the WordPress baseline. Task completion rose from 30\% to 95\%. Median creation time dropped from 45 minutes to 3.2 minutes. The System Usability Scale score improved from 38.0 to 82.5---moving from ``Poor'' to ``Excellent'' on Brooke's scale. The results suggest that pairing user-centred design with generative AI can meaningfully lower the barrier to creating interactive pedagogical content for university teaching.

\bigskip
\noindent
\textbf{Keywords:} H5P, interactive video, educational technology, AI-assisted content authoring, usability evaluation, user-centred design.
